% Capitolo 5 - Conclusioni
% Ultima modifica: 2026-05-29
% Fonti principali usate: MEF_RGS2025, MEF_RGS2025_NdA, Bosi2018, Fornero2018,
%   FrancoTommasino2020, OCPI2026, AcemogluRestrepo2020, AcemogluJohnson2023,
%   Bacchiocchi2025, Artoni2014, Padula2024

%% ═══════════════════════════════════════════════════════════════════
\chapter*{Conclusioni}
\addcontentsline{toc}{chapter}{Conclusioni}
%% ═══════════════════════════════════════════════════════════════════

\noindent
Questa tesi è partita da una domanda che il dibattito previdenziale italiano
tende a non porre: non quando si va in pensione o quanto costa il sistema, ma
se la pensione che attende chi entra oggi nel mercato del lavoro sarà
sufficiente. La risposta costruita nei quattro capitoli è netta nei due piani
che la compongono e scomoda nella loro combinazione. Sul piano contabile il
sistema regge: la spesa pensionistica tocca il picco del 17,1\% del PIL intorno
al 2042-2043 e rientra al 14,0\% nel 2070 grazie ai meccanismi automatici del
metodo contributivo \cite{MEF_RGS2025_NdA}. Sul piano dell'adeguatezza lo stesso
sistema arretra: il tasso di sostituzione lordo del dipendente privato scende dal
73,6\% del 2010 al 58,4\% del 2070 \cite{MEF_RGS2025}, e per le carriere
discontinue tipiche della generazione che entra oggi il valore effettivo è più
basso ancora. La riforma del 1995 non ha eliminato il rischio previdenziale: lo
ha spostato, dalla collettività che lo assorbiva con il metodo retributivo al
singolo lavoratore che oggi lo porta in dote.

Le leve per ricostruire adeguatezza esistono, e il terzo capitolo le ha misurate:
pensione contributiva di garanzia, perequazione selettiva sulle pensioni alte,
separazione contabile tra previdenza e assistenza, rafforzamento del secondo
pilastro come add-on volontario. Nessuna è bloccata da un limite tecnico, ma
dall'economia politica di un paese in cui l'elettore decisivo ha già maturato i
propri diritti, e le riforme che ripristinano equità intergenerazionale chiedono un
sacrificio a chi vota oggi per un beneficio a chi voterà fra trent'anni. Che il
vincolo sia politico e non contabile lo dimostra il costo di smontare gli
automatismi già in vigore: bloccare l'adeguamento dei requisiti e dei coefficienti
aggiungerebbe circa cinquantotto punti di debito sul PIL al 2070
\cite{MEF_RGS2025_NdA}. La conclusione del terzo capitolo non è un programma di
riforma ma una gerarchia di fattibilità, dove le riforme efficienti sono
politicamente improbabili e quelle disponibili sono parziali.

Si potrebbe obiettare che basterebbe riformare il pilastro pubblico, senza chiedere
nulla all'individuo. La gerarchia di fattibilità risponde: quelle riforme costano
punti di PIL che l'economia politica non concede, ed è questo a rendere l'azione
individuale necessaria e non sostitutiva. Resta la leva che non dipende dal
legislatore. Il quarto capitolo ha mostrato che un lavoratore informato dispone di
strumenti concreti e immediati, dal conferimento del TFR a un fondo a basso costo
alla deduzione fiscale fino a 5.164,57 euro che trasforma il versamento in
rendimento garantito prima dell'investimento, fino alla linea dinamica e al piano di
accumulo su ETF. Per chi ha reddito e carriera continua, come il Profilo~A, questi
strumenti recuperano sette-otto punti di tasso di sostituzione e agiscono subito,
senza attendere alcuna legge: per questo segmento l'azione individuale è la leva
decisiva, e la tesi la difende come tale.

Lo stesso esercizio, però, non si limita a prescrivere. Misura anche dove la
prescrizione non arriva. I profili B, C e D mostrano che la capacità di esercitare
quelle mosse dipende dal reddito disponibile e dalla continuità contributiva, e a
queste due disuguaglianze se ne aggiunge una terza, l'alfabetizzazione finanziaria,
che la letteratura documenta più bassa fra i giovani e fra le donne e crescente
con il reddito e l'istruzione \cite{LusardiMitchell2014}. I
tre svantaggi si cumulano sulle stesse persone. Questo non rende il quarto capitolo
irrilevante per i profili vulnerabili: ne è semmai la parte più utile, perché
quantifica il divario residuo che l'azione individuale non chiude, ed è la misura
di quel divario a dimensionare il trasferimento pubblico, la pensione contributiva
di garanzia del terzo capitolo, che è l'unico strumento in grado di colmarlo. I
due pilastri non competono per la stessa platea: l'add-on volontario lavora sopra
la soglia di adeguatezza, per chi ha il margine di salirci; il pavimento garantito
lavora sotto, per chi quel margine non ce l'ha. L'azione individuale è la risposta
più rapida e a costo politico nullo, e va difesa e diffusa. Non è, da sola, una
risposta sistemica, perché la quota di lavoratori che ne trae più beneficio non
coincide con quella che ne avrebbe più bisogno: ed è per questo, non malgrado
questo, che il pavimento pubblico resta necessario.

Qui si chiude il paradosso da cui la tesi è partita e si apre quello che lascia
in eredità. Un sistema progettato per garantire un tenore di vita dignitoso nella
vecchiaia potrebbe non garantirlo a una quota larga della generazione che entra
oggi. Non è una previsione catastrofista: è l'algebra del metodo contributivo
applicata a un quarto di secolo di crescita reale italiana attestata attorno allo
0,5\% medio annuo \cite{MEF_RGS2025}. La sostenibilità finanziaria resta una
condizione necessaria ma non sufficiente: i conti tornano, eppure per molti le
pensioni potrebbero non bastare.

\section*{La responsabilità individuale come principio e come alibi}
\addcontentsline{toc}{section}{La responsabilità individuale come principio e come alibi}

\noindent
Una domanda normativa attraversa l'intera tesi senza sciogliersi: è giusto che la
responsabilità di una pensione adeguata ricada sul singolo lavoratore? In un quadro
liberale coerente la risposta è affermativa, perché lo Stato garantisce la soglia
minima e ciò che la supera è frutto dell'agentività individuale, e la neutralità
attuariale, a parità di contributi parità di pensione, è prima di tutto un principio
di giustizia e non un tecnicismo \cite{Bosi2018}. La risposta regge, però, solo a tre
condizioni che in Italia oggi non si verificano: l'informazione, l'accessibilità
degli strumenti e la crescita.

L'informazione viene prima di tutto, perché chiedere una scelta previdenziale a chi
non conosce il proprio tasso di sostituzione atteso non è libertà ma delega in
bianco, e la comunicazione individuale prevista già dalla riforma Dini non è mai
stata resa sistematica, lasciando la maggioranza dei lavoratori all'oscuro della
pensione che li attende \cite{Fornero2018, Padula2024}. Segue l'accessibilità degli
strumenti, perché la deduzione fiscale è una leva potente ma il mercato dei fondi non
è uniforme, e i piani individuali distribuiti dalle reti commerciali presentano costi
tra i più alti d'Europa e un'asset allocation troppo conservativa per orizzonti
trentennali \cite{FrancoTommasino2020}: è qui che lo Stato arbitro, non gestore, ha
un compito preciso, da assolvere con la trasparenza obbligatoria e con la
concorrenza, non con l'obbligo di adesione. Resta la crescita, perché la logica
attuariale produce risultati equi solo se ha remunerato i contributi versati, e con
un PIL nominale cresciuto in termini reali dello 0,5\% medio annuo in un quarto di
secolo \cite{MEF_RGS2025} il rendimento implicito del sistema è rimasto sotto
qualunque alternativa di risparmio ragionevolmente accessibile.

Riconoscerlo non significa abbandonare la responsabilità individuale, ma chiedere
allo Stato ciò che in un quadro liberale gli compete prima che la partita inizi, cioè
garantire regole uguali, chiare e accessibili a tutti: informazione obbligatoria,
concorrenza sui costi dei fondi e separazione netta tra previdenza e assistenza non
intaccano la neutralità attuariale, anzi chiedono che valga per tutti e non soltanto
per chi partiva già nelle condizioni di sfruttarla \cite{Artoni2014}. Il sistema
contributivo è tecnicamente corretto e le proiezioni della Ragioneria sono credibili,
eppure un sistema che trasferisce sull'individuo il rischio demografico e di
longevità senza garantirgli informazione, strumenti e crescita per gestirlo non è
ancora un sistema equo. La tensione tra rigore contabile e giustizia sociale che
questa tesi ha esaminato non si scioglie in queste pagine: è il problema che la
generazione entrata oggi nel mercato del lavoro eredita, insieme ai contributi da
versare.
